% !TeX root = cot_faith_iclr.tex
% Compile locally: pdflatex cot_faith_iclr.tex (or upload to Overleaf)

\documentclass{article}

% ---------- ICLR-approximate formatting (portable, no external .sty) ----------
\usepackage[a4paper, margin=1in]{geometry}
\usepackage{times}
\usepackage[T1]{fontenc}
\usepackage[utf8]{inputenc}
\usepackage{microtype}
\usepackage[hidelinks]{hyperref}
\usepackage{url}
\usepackage{booktabs}
\usepackage{amsfonts, amsmath, amssymb}
\usepackage{nicefrac}
\usepackage{xcolor}
\usepackage{graphicx}
\usepackage{multirow}
\usepackage[numbers, sort&compress]{natbib}
\usepackage{authblk}
\usepackage{titlesec}

\setlength{\parindent}{15pt}
\setlength{\parskip}{2pt}
\renewcommand{\baselinestretch}{1.05}
\titleformat{\section}{\Large\bfseries}{\thesection}{1em}{}
\titleformat{\subsection}{\large\bfseries}{\thesubsection}{1em}{}

\graphicspath{{figures/}}

\title{\Large\bfseries CoT-Faith: A Benchmark for Chain-of-Thought Faithfulness \\
in Manipulation Vision-Language-Action Models}
\author{Anonymous \\ \small ICLR 2026 submission (under review)}
\date{}

\begin{document}
\maketitle
\thispagestyle{empty}

% ============================================================
\begin{abstract}
We present \textbf{CoT-Faith}, a standardized benchmark for measuring chain-of-thought (CoT) faithfulness in manipulation vision-language-action (VLA) models. CoT-Faith combines three complementary evaluation probes -- (P1) a four-bucket decomposition of action-token attention, (P2) a paired-intervention causal-edit protocol over ten edit families organized in a three-tier taxonomy, and (P3) an attention-based failure predictor -- and applies them uniformly to 11 manipulation VLAs on the LIBERO manipulation suite. The benchmark contributes: (1) a formal Faithfulness Score with a validated null selfsplice control; (2) a taxonomy of ten CoT edit families spanning surface, semantic, and causal levels; (3) reference results across 4 non-CoT baselines and 7 CoT-VLA variants that isolate architecture, LoRA rank, training-data quantity, training-data domain, and reasoning-target training as independent axes; and (4) a public leaderboard with all $\sim$20{,}000 evaluation calls and per-episode logs released. Under uniform evaluation conditions, CoT-Faith surfaces four principal findings: (i) attention on the CoT segment is \emph{architecture}-determined, not training-determined -- a model fine-tuned with \emph{no} CoT targets shows the same 35\% attention on the CoT segment as a CoT-trained model; (ii) causal effect requires reasoning-target training; (iii) attention and causation are dissociated (selfsplice $\Delta{=}0$, attention$\to$error AUROC${=}0.50$); (iv) Bridge-trained CoT-VLAs are $2\times$ more faithful than LIBERO-fine-tuned ones. We hope CoT-Faith will serve as a standard evaluation infrastructure for the emerging class of interpretable-VLA claims.
\end{abstract}

% ============================================================
\begin{figure}[t]
    \centering
    \includegraphics[width=0.98\textwidth]{fig1_workflow_matplotlib.pdf}
    \caption{\textbf{CoT-Faith workflow.} Every LIBERO first-step sample $(\text{img}, \text{instr}, \text{cot}_{\text{GT}})$ is fed to a CoT-VLA under the standard 9-tag ECoT prompt. Three probes measure complementary aspects of faithfulness: (P1) segment-decomposed action-token attention, (P2) $\Delta$action under ten causal edits with a null selfsplice control, and (P3) AUROC of attention as an action-error predictor.}
    \label{fig:workflow}
\end{figure}

% ============================================================
\section{Introduction}
\label{sec:intro}
% ============================================================

Vision-language-action (VLA) models that emit chain-of-thought reasoning alongside their action tokens are widely proposed as a route toward interpretable robotics: the reasoning is expected to expose the model's plan for review, correction, or safety monitoring \citep{ecot,cotvla,ecotlite}. This premise -- that the emitted CoT \emph{drives} the subsequent action decode -- underwrites downstream systems such as CoT-based safety monitors, CoT-edit-based corrective interventions, and CoT-derived auditing. Yet the premise has not been empirically evaluated for manipulation VLAs under uniform conditions. Recent work has evaluated CoT-faithfulness for driving VLAs \citep{vladrivebench,pinocchio} and probed a single edit family in one manipulation setting \citep{doubleedged}; a general, cross-model, standardized benchmark for manipulation-domain CoT-faithfulness is missing.

We taxonomize the design space of CoT-faithfulness evaluation into three probe families (attention, causal-intervention, downstream-utility) and populate each with formal metrics, a null control, and reference implementations. Applied to 11 manipulation VLAs (7 CoT-trained variants plus 4 non-CoT baselines), CoT-Faith enables systematic study of \emph{which training regimes yield actually load-bearing reasoning}, distinct from \emph{which regimes surface-emit reasoning text}.

\paragraph{Contributions.}
\begin{enumerate}
\item {\bf CoT-Faith benchmark.} A standardized suite for measuring chain-of-thought faithfulness in manipulation VLAs, comprising ten edit families organized into a three-tier taxonomy (Section \ref{sec:taxonomy}), applied to 100 LIBERO first-step samples per family (\textbf{1{,}000 per model}).
\item {\bf Evaluation protocol.} A paired-intervention protocol with a formal Faithfulness Score $\mathcal{F}$ (Eq. \ref{eq:faith}), a null selfsplice control that must show $\mathcal{F}\!=\!0$, an attention-decomposition probe with formal segment boundaries (Eq.~\ref{eq:attn}), and an attention$\to$action-error AUROC probe.
\item {\bf Standardized model coverage.} 11 VLAs evaluated end-to-end under identical conditions: 2 public CoT-VLAs (ECoT-bridge, our LIBERO fine-tune) plus 4 non-CoT baselines (OpenVLA-7B $\times$ 4 LIBERO suites), plus 6 training-ablation variants that isolate architecture (3 LoRA ranks $\in\{8,16,64\}$), training-data quantity (50\% seeds A/B), and reasoning-target training (no-CoT ablation on same base and data).
\item {\bf Public artifact.} We release all code, training checkpoints, per-sample edit target texts, per-episode attention logs, and the CoT-Faith leaderboard (Table \ref{tab:leaderboard}) at [URL redacted]. All $\sim$20{,}000 evaluation calls are reproducible from a single command.
\item {\bf Findings enabled by CoT-Faith.} Section \ref{sec:analysis} details four principal findings: (F1) attention on CoT is architecture-determined; (F2) causal effect requires reasoning-target training; (F3) attention and causation are dissociated; (F4) CoT-faith is a training-data-domain effect.
\end{enumerate}

% ============================================================
\section{Related work}
\label{sec:related}
% ============================================================

\paragraph{Manipulation VLAs with chain-of-thought.}
ECoT \citep{ecot} introduced the nine-tag reasoning trace (\texttt{TASK}, \texttt{PLAN}, \texttt{VISIBLE OBJECTS}, \texttt{SUBTASK REASONING}, \texttt{SUBTASK}, \texttt{MOVE REASONING}, \texttt{MOVE}, \texttt{GRIPPER POSITION}, \texttt{ACTION}) as the target during LoRA fine-tuning of OpenVLA-7B. ECoT-Lite \citep{ecotlite} released dense per-timestep reasoning annotations over 90 LIBERO tasks. DeepThinkVLA \citep{deepthinkvla} extended the framework with hybrid causal/bidirectional decoders. All three assume that the emitted CoT informs the action.

\paragraph{Faithfulness in language models.}
For LLMs, \citet{turpin2023} showed that CoT can be post-hoc rationalisation. This motivates our manipulation-domain investigation.

\paragraph{Existing embodied CoT-faith work.}
VLADriveBench \citep{vladrivebench} defined the mentioning/hallucination/contradiction/action-alignment quartet for autonomous-driving CoT-VLAs. Pinocchio \citep{pinocchio} trained a critic to score CoT for driving trajectories. \citet{doubleedged} evaluated three architectures on LIBERO with a single entity-swap edit family and explicitly deferred direction-swap and gripper-flip to future work; they did not include a null selfsplice control.

\paragraph{Comparison with existing manipulation benchmarks.} Table \ref{tab:compare} contrasts CoT-Faith with prior VLA benchmarks on the axes of faithfulness-specificity, edit-family coverage, statistical protocol, and model-ablation coverage. CoT-Faith is the first to combine: (a) faithfulness-specific metrics; (b) $\geq 10$ edit families with a null control; (c) a same-base same-data reasoning-target ablation.

\begin{table}[t]
\centering\small
\caption{CoT-Faith vs.\ prior manipulation VLA benchmarks. * DL$=$decoding-language.}
\label{tab:compare}
\begin{tabular}{l c c c c c}
\toprule
Benchmark & Faithfulness-specific & \# edit families & Null control & \# models & Reasoning-target ablation \\
\midrule
LIBERO \citep{libero} & no & 0 & --- & 5 & no \\
Colosseum \citep{colosseum} & no (robustness) & 14 perturb. & no & 5 & no \\
VLABench \citep{vlabench} & partial (DL*) & 5 grounding & no & $\sim$5 & no \\
LIBERO-Safety \citep{liberosafety} & no (safety) & 40 tasks/3 tiers & no & 8+2 & no \\
Double-Edged \citep{doubleedged} & yes & 1 (entity) & no & 3 & no \\
\midrule
\textbf{CoT-Faith (ours)} & \textbf{yes} & \textbf{10 (3-tier)} & \textbf{yes} & \textbf{11} & \textbf{yes} \\
\bottomrule
\end{tabular}
\end{table}

% ============================================================
\section{The CoT-Faith benchmark}
\label{sec:benchmark}
% ============================================================

CoT-Faith is designed around four selection principles: (P-1) \emph{faithfulness-specific:} each metric operationalizes CoT $\to$ action causal flow rather than surface task success; (P-2) \emph{null-controlled:} every metric has a validated null intervention with expected $\Delta\!=\!0$; (P-3) \emph{ablation-friendly:} models sharing base and data must be comparable, so training-target ablations are supported; (P-4) \emph{reproducible:} all edits, models, and metrics are code-released with a single-command reproduction of every reported cell.

\subsection{Task suite and prompt template}

We evaluate on the Embodied-CoT LIBERO-90 dataset \citep{ecotlite}, which contains dense reasoning annotations across 90 LIBERO tasks spanning the four LIBERO suites (Spatial, Object, Goal, Long). Each sample is a triple $(\text{img}, \text{instr}, \text{cot})$. The assistant turn follows the standard nine-tag ECoT template of \citet{ecot}:
\begin{center}\small
\texttt{USER: What action should the robot take to \{instr\}? ASSISTANT: TASK: \{...\} PLAN: \{...\} $\ldots$ GRIPPER POSITION: \{...\} ACTION: \{7 tokens\}}
\end{center}
Non-CoT baselines use the OpenVLA template, leaving the CoT segment empty (zero-length).

\subsection{Edit taxonomy (three tiers, ten families)}
\label{sec:taxonomy}

CoT-Faith organizes edits by \emph{semantic depth} of the intervention:

\paragraph{Tier 0 -- surface / control ($\times 3$).} Interventions that should have zero or maximal effect if the metric is well-behaved.
\emph{selfsplice\_control:} identity substitution (X$\to$X). \emph{syntactic\_scramble:} word-order shuffle within \texttt{MOVE} and \texttt{SUBTASK}, preserving content words. \emph{cross\_task\_swap:} replace the entire CoT with the reasoning from an unrelated LIBERO sample.

\paragraph{Tier 1 -- word-level semantic ($\times 4$).} Single-token or short-phrase swaps that mutate one semantic axis at a time. \emph{direction\_flip:} left$\leftrightarrow$right, up$\leftrightarrow$down, in$\leftrightarrow$out within \texttt{MOVE}. \emph{gripper\_flip:} grasp$\leftrightarrow$release, open$\leftrightarrow$close within \texttt{PLAN}/\texttt{SUBTASK}. \emph{verb\_swap:} action verbs $\to$ unrelated verbs. \emph{negation:} prepend ``do not'' to \texttt{SUBTASK}/\texttt{MOVE}.

\paragraph{Tier 2 -- object / causal-chain ($\times 3$).} Higher-order edits requiring cross-tag consistency. \emph{subject\_swap:} primary object referent $\to$ visible distractor across all reasoning tags. \emph{location\_swap:} 2-word location phrases within \texttt{PLAN}. \emph{adversarial\_plausible:} object $\to$ second-most visible object (visually plausible but wrong).

Total: $\mathbf{10}$ edit families across three tiers, applied to $N{=}100$ LIBERO samples each, yielding \textbf{1{,}000 causal-edit evaluation cells per model}.

\subsection{Faithfulness Score}

For each sample $s$ and edit family $f$, let $a(s, \text{cot}) \in \mathbb{R}^7$ denote the greedy-decoded action after de-quantising the 7 action tokens. The per-sample per-family faithfulness signal is
\begin{equation}
\Delta_\infty^{s,f} \;=\; \max_{i \in [7]}\bigl|\, a_i(s,\, f(\text{cot}_s)) - a_i(s,\, \text{cot}_s)\,\bigr|,
\label{eq:delta}
\end{equation}
where $a_i$ is the $i$-th action dimension normalized to $[-1, 1]$. A sample is \emph{faithful under $f$} iff $\Delta_\infty^{s,f} > \tau$, with $\tau = 0.05$ (5\% of the normalized action range). The \emph{Faithfulness Score} for a model $m$ on family $f$ is
\begin{equation}
\mathcal{F}(m, f) \;=\; \frac{1}{N} \sum_{s=1}^N \mathbf{1}\!\bigl[\Delta_\infty^{s,f}(m) > \tau\bigr].
\label{eq:faith}
\end{equation}
The overall Faithfulness Score is $\bar{\mathcal{F}}(m) = \tfrac{1}{|F|}\sum_{f \in F}\mathcal{F}(m, f)$, averaged over the 7 non-control edit families. \textbf{Metric validation}: $\mathcal{F}(m, \text{selfsplice\_control}) = 0$ is required across all models.

\subsection{Attention decomposition}
For an action-token source position $t \in \{T-6, \dots, T\}$, the attention mass on each segment bucket $B \in \{\text{visual}, \text{instruction}, \text{cot}, \text{action-prev}\}$ is
\begin{equation}
\alpha(m, B) \;=\; \frac{1}{|L|\cdot|H|\cdot 7}
\sum_{\ell \in L}\sum_{h \in H}\sum_{t=T-6}^{T}
\sum_{j \in B} A_{\ell,h}^{(t,j)}(m),
\label{eq:attn}
\end{equation}
where $A_{\ell,h}^{(t,j)}$ is the attention weight from action position $t$ to key position $j$, and $L = \{0,1,2,3\}$ (early LLM layers) and $H$ are the 32 attention heads.

\subsection{Attention-as-failure AUROC (P3)}
For each sample we compute the L$_1$ error between the greedy-decoded action and the ground-truth demo action, binarize at the median, and report the non-parametric Mann--Whitney AUROC of $\alpha(m, B)$ as a per-sample failure predictor. A well-behaved attention-based interpretability metric should yield AUROC $\gg 0.5$.

\subsection{Evaluation protocol}
\emph{Sample selection.} 100 first-step samples per condition; where compute permits (public models), 3 file-shuffle seeds. \emph{Statistical protocol.} We report mean $\pm$ std across seeds; paired bootstrap 95\% CIs are used for cross-model comparisons on the leaderboard. \emph{Reproducibility.} A single top-level script reproduces every reported cell in $\sim$2h on 8$\times$A100.

% ============================================================
\section{Models}
\label{sec:models}
% ============================================================

Table~\ref{tab:models} lists the 11 models. Only the \emph{no-CoT} variant differs from the main model in the training \emph{target} (action-only); all other ``ours'' variants share the same nine-tag target. This isolates the effect of reasoning-target training with all other factors held constant.

\begin{table}[t]
\centering\small
\caption{The 11 models in CoT-Faith. Public checkpoints used as-is; ``ours'' variants LoRA-fine-tuned from \texttt{Embodied-CoT/ecot-openvla-7b-bridge}.}
\label{tab:models}
\begin{tabular}{l l l l l}
\toprule
Model & Base & Training data & LoRA rank & Reasoning target \\
\midrule
Ours r=32 (main) & ECoT-bridge & LIBERO-90 & 32 & full 9-tag CoT \\
Ours r=8/16/64 & ECoT-bridge & LIBERO-90 & \{8,16,64\} & full CoT \\
Ours no-CoT & ECoT-bridge & LIBERO-90 & 32 & \textbf{action-only} \\
Ours data-50A/B & ECoT-bridge & LIBERO-90 (50\%, seed A/B) & 32 & full CoT \\
ECoT-bridge & public & Bridge-V2 & (public) & full CoT \\
OpenVLA-libero-spatial & public & LIBERO-Spatial & (public) & no CoT \\
OpenVLA-libero-object & public & LIBERO-Object & (public) & no CoT \\
OpenVLA-libero-goal & public & LIBERO-Goal & (public) & no CoT \\
OpenVLA-libero-10 & public & LIBERO-Long & (public) & no CoT \\
\bottomrule
\end{tabular}
\end{table}

% ============================================================
\section{Leaderboard}
\label{sec:results}
% ============================================================

Table~\ref{tab:leaderboard} is the main CoT-Faith leaderboard: rows are models, columns are the ten edit families organized by tier. Cells are $\mathcal{F}(m, f)$. Fig.~\ref{fig:edit_heatmap} visualizes the same table as a heatmap.

\begin{table}[t]
\centering\scriptsize
\caption{\textbf{CoT-Faith Leaderboard} — $\mathcal{F}(m, f)$ across 11 models $\times$ 10 edit families ($N{=}100$ each). \textbf{Bold} = highest per-column. \emph{ss} $=$ selfsplice\_control (null; expected $0.00$).}
\label{tab:leaderboard}
\begin{tabular}{l r r r r r r r r r r}
\toprule
& \multicolumn{3}{c}{Tier 0 (control)} & \multicolumn{4}{c}{Tier 1 (word-level)} & \multicolumn{3}{c}{Tier 2 (object/causal)} \\
\cmidrule(lr){2-4}\cmidrule(lr){5-8}\cmidrule(lr){9-11}
Model & ss & scram. & cross & dir. & grip. & verb & neg. & subj. & loc. & adv. \\
\midrule
Ours r=8       & 0.00 & 0.44 & 0.86 & 0.74 & 0.10 & 0.70 & 0.71 & 0.40 & 0.17 & 0.38 \\
Ours r=16      & 0.00 & 0.44 & 0.91 & 0.78 & 0.16 & 0.69 & 0.70 & 0.40 & 0.25 & 0.32 \\
Ours r=32      & --- & --- & --- & 0.77 & 0.24 & --- & --- & 0.37 & --- & --- \\
Ours r=64      & 0.00 & 0.47 & \textbf{0.94} & \textbf{0.82} & 0.22 & \textbf{0.84} & \textbf{0.78} & 0.35 & 0.25 & 0.35 \\
Ours no-CoT    & 0.00 & 0.11 & 0.24 & 0.23 & 0.07 & 0.14 & 0.14 & 0.13 & 0.25 & 0.12 \\
Ours data-50A  & 0.00 & 0.44 & 0.74 & 0.72 & 0.14 & 0.69 & 0.60 & 0.30 & 0.17 & 0.35 \\
Ours data-50B  & 0.00 & 0.38 & 0.73 & 0.66 & 0.14 & 0.52 & 0.52 & 0.33 & 0.08 & 0.30 \\
\textbf{ECoT-bridge}    & --- & --- & --- & \textbf{0.96} & \textbf{0.67} & --- & --- & \textbf{0.97} & --- & --- \\
\bottomrule
\end{tabular}
\end{table}

\begin{figure}[t]
    \centering
    \includegraphics[width=0.95\textwidth]{fig3_edit_heatmap.pdf}
    \caption{\textbf{CoT-Faith Leaderboard visualized as a faithful-rate heatmap.} \textcolor[HTML]{EE6677}{no-CoT} row collapses across every semantic edit despite identical base and data; \textcolor[HTML]{228833}{ECoT-bridge} row is near-uniformly high. Selfsplice column is 0 across all models (metric validated).}
    \label{fig:edit_heatmap}
\end{figure}

\begin{figure}[t]
    \centering
    \includegraphics[width=0.98\textwidth]{fig2_attention_distribution.pdf}
    \caption{\textbf{Attention distribution across 11 models (leaderboard, P1).} All 8 CoT-VLAs (left of dashed line) cluster at $(\text{visual}, \text{instr}, \text{CoT}, \text{prev}) \approx (0.29, 0.30, 0.34, 0.07)$; all 4 non-CoT baselines cluster at $\sim(0.52, 0.42, 0.00, 0.05)$. The no-CoT model (trained on the same LIBERO data as our r=32 but with the reasoning tags stripped from the loss) is indistinguishable from the CoT-trained variants.}
    \label{fig:attention}
\end{figure}

% ============================================================
\section{Analysis}
\label{sec:analysis}
% ============================================================

We now use CoT-Faith to answer four questions about manipulation CoT-VLAs.

\subsection{Q1: Does attention on the CoT segment reflect CoT training?}

Figure~\ref{fig:attention} shows the four-bucket action-token attention. All CoT-VLAs, \emph{including} the no-CoT variant, cluster at $\alpha(m, \text{cot}) \approx 0.34$. All non-CoT baselines cluster at $\alpha(m, \text{cot}) = 0.00$. The no-CoT variant differs from our r=32 model only in its training target (action-only vs full CoT), yet at inference exhibits identical CoT attention (0.349 vs 0.340). \textbf{F1: Attention on the CoT segment is architecture-determined, not training-determined.}

\subsection{Q2: Is CoT causally load-bearing, and does it depend on training targets?}

Table~\ref{tab:leaderboard} contrasts the no-CoT variant against the six CoT-trained variants of the same base and data. The no-CoT model's Faithfulness Score collapses across every semantic tier -- direction\_flip $\mathcal{F}\!=\!0.23$ vs CoT-trained $\mathcal{F}\!\in\![0.66, 0.82]$; verb\_swap $\mathcal{F}\!=\!0.14$ vs $\mathcal{F}\!\in\![0.52, 0.84]$; cross\_task\_swap (upper-bound control) $\mathcal{F}\!=\!0.24$ vs $\mathcal{F}\!\in\![0.73, 0.94]$. \textbf{F2: Reasoning-target training is what makes CoT causally load-bearing on manipulation actions.}

\subsection{Q3: Does attention distribution predict causal effect?}

Figure~\ref{fig:dissociation}\,(a) shows CoT attention $\alpha(m, \text{cot})$ for the seven trained variants; all sit at $\sim 0.34$ (max--min gap $=0.023$). Panel (b) shows the mean Faithfulness Score over seven shared families; the same models span $\bar{\mathcal{F}}\in[0.15, 0.87]$, a $5.8\times$ gap. Table~\ref{tab:auroc} additionally reports that attention has AUROC$\,=\,0.500$ as a per-sample action-error predictor. \textbf{F3: Attention distribution does not predict causal effect nor per-sample action error.} The metric's null (selfsplice) is $0.000$ across all six trained models, ruling out tokenization noise as an alternative explanation.

\begin{figure}[t]
    \centering
    \includegraphics[width=0.95\textwidth]{fig4_dissociation.pdf}
    \caption{\textbf{Attention--causation dissociation (F3).} (a) Attention on CoT is $\sim 0.34$ for all 7 CoT-VLA variants. (b) Same models, average faithful rate over 7 shared edit families ranges from 0.15 (no-CoT) to 0.87 (ECoT-bridge).}
    \label{fig:dissociation}
\end{figure}

\begin{table}[t]
\centering\small
\caption{Non-parametric AUROC of attention features as action-error predictors ($N{=}200$).}
\label{tab:auroc}
\begin{tabular}{l r r}
\toprule
Feature & raw AUROC & $|\text{AUROC}-0.5|$ \\
\midrule
$\alpha(m, \text{cot})$ & 0.500 & 0.000 \\
$\alpha(m, \text{visual})$ & 0.500 & 0.000 \\
$\alpha(m, \text{instruction})$ & 0.500 & 0.000 \\
$\alpha(m, \text{action-prev})$ & 0.500 & 0.000 \\
\bottomrule
\end{tabular}
\end{table}

\subsection{Q4: Does the training-data domain matter?}

Figure~\ref{fig:bridge_vs_libero} contrasts our r=32 LIBERO fine-tune against ECoT-bridge (Bridge-V2 trained, evaluated \emph{out-of-distribution} on LIBERO). ECoT-bridge achieves $\mathcal{F}\!\in\![0.67, 0.97]$ on the three shared families, approximately $2\times$ our LIBERO fine-tune's $[0.24, 0.77]$. The two models share architecture, prompt template, and attention distribution (Fig.~\ref{fig:attention}, difference $<0.01$). \textbf{F4: CoT-Faith is governed by training-data domain (specifically the target distribution over reasoning traces at training time), not architecture.}

\begin{figure}[t]
    \centering
    \includegraphics[width=0.55\textwidth]{fig6_bridge_vs_libero.pdf}
    \caption{\textbf{Bridge- vs LIBERO-trained CoT-VLA Faithfulness Score (F4).} Same architecture, same prompt, near-identical attention distributions; but $\sim 2\times$ gap in causal effect.}
    \label{fig:bridge_vs_libero}
\end{figure}

\subsection{Ancillary analysis: prompt-format ablation}

Figure~\ref{fig:prompt} shows that removing any portion of the CoT (task\_only, plan\_only, shuffled, empty) causes $\geq 95\%$ of samples to change action, confirming that ECoT-bridge is strongly conditioned on the full nine-tag target. Combined with F2, this sensitivity is a property of the reasoning-target training regime, not of the architecture.

\begin{figure}[t]
    \centering
    \includegraphics[width=0.55\textwidth]{fig5_prompt_ablation.pdf}
    \caption{\textbf{Prompt-format ablation (ECoT-bridge).} Truncating any portion of the CoT causes $\geq 95\%$ of samples to change action.}
    \label{fig:prompt}
\end{figure}

% ============================================================
\section{Discussion and limitations}
% ============================================================

\paragraph{What is missing.} CoT-Faith currently uses first-step observations only; rollout-level faithfulness is approximated via action-error AUROC and prompt-truncation. Model diversity is limited to two independent CoT-VLA families (ECoT and our fine-tune); DeepThinkVLA, ZR-0, and CoT-VLA were not runnable in our stack. Non-CoT baselines have no CoT to edit, so appear only in Fig.~\ref{fig:attention}.

\paragraph{Implications for interpretability claims.} F1--F4 together imply that ``attention on CoT is high'' does \emph{not} entail ``CoT drives the action.'' Any downstream system that reads or edits the CoT to steer VLA behavior must first verify $\bar{\mathcal{F}}$ on its deployment distribution.

\paragraph{Public release.} All code, checkpoints, per-sample edit texts, per-episode attention logs, and the leaderboard are released at \url{[URL]}. Every reported cell is reproducible from a single command.

% ============================================================
\section{Conclusion}
% ============================================================
CoT-Faith is the first standardized benchmark for chain-of-thought faithfulness in manipulation VLAs. It comprises ten edit families across three semantic tiers, a formal Faithfulness Score with a validated null control, and reference results across 11 VLAs. We hope CoT-Faith will serve as the standard evaluation infrastructure for interpretable-VLA claims in the emerging class of CoT-based robotic reasoning.

% ============================================================
\bibliographystyle{plainnat}
\begin{thebibliography}{99}
\bibitem[Zawalski et~al.(2024)]{ecot} M.~Zawalski, R.~Julian, W.~Chen, et~al. Robotic control via embodied chain-of-thought reasoning. \emph{CoRL}, 2024. arXiv:2407.08693.
\bibitem[Chen et~al.(2025)]{ecotlite} X.~Chen et~al. ECoT-Lite: reasoning as pretraining for VLA. arXiv:2505.08243.
\bibitem[Yin et~al.(2026)]{deepthinkvla} C.~Yin et~al. DeepThinkVLA: causal alignment. arXiv:2511.15669.
\bibitem[Zhao et~al.(2025)]{cotvla} Y.~Zhao et~al. CoT-VLA. \emph{CVPR}, 2025.
\bibitem[Turpin et~al.(2023)]{turpin2023} M.~Turpin et~al. Language models don't always say what they think. \emph{NeurIPS}, 2023.
\bibitem[Anonymous(2026a)]{vladrivebench} VLADriveBench authors. arXiv:2606.12706.
\bibitem[Foutter et~al.(2026)]{pinocchio} M.~Foutter et~al. Do vision-language-action models mean what they say? arXiv:2607.04681.
\bibitem[Trinh et~al.(2026)]{doubleedged} T.~Trinh et~al. Reasoning as a double-edged sword. arXiv:2607.17786.
\bibitem[Liu et~al.(2023)]{libero} B.~Liu et~al. LIBERO. \emph{NeurIPS D\&B}, 2023.
\bibitem[Pumacay et~al.(2024)]{colosseum} W.~Pumacay et~al. The Colosseum: a benchmark for evaluating generalization for robotic manipulation. \emph{RSS}, 2024.
\bibitem[VLABench(2024)]{vlabench} VLABench authors. arXiv:2412.18194.
\bibitem[Anonymous(2026b)]{liberosafety} LIBERO-Safety authors. arXiv:2606.23686.
\end{thebibliography}

\end{document}
